% Präambel
\documentclass[11pt,a4paper,oneside, 
liststotoc, 					% Tabellen- und Abbildungsverzeichnis ins Inhaltsverzeichnis
bibtotoc,						% Literaturverzeichnis ins Inhaltsverzeichnis aufnehmen
titlepage, 						% Titlepage-Umgebung statt \maketitle
headsepline, 					% horizontale Linie unter Kolumnentitel
%abstracton,					% Überschrift beim Abstract einschalten, Abstract muss dazu in {abstract}-Umgebung stehen
%DIV11,							% auskommentieren, um den Seitenspiegel zu vergrößern
BCOR6mm,						% Bindekorrektur, die den Seitenspiegel um 6mm nach rechts verschiebt,
]{scrreprt}			
\usepackage{ucs} 				% Dokument in utf8-Codierung schreiben und speichern
\usepackage[utf8x]{inputenc} 	% ermöglicht die direkte Eingabe von Umlauten
\usepackage[ngerman]{babel} 	% deutsche Trennungsregeln und Übersetzung der festcodierten Überschriften
\usepackage[T1]{fontenc} 		% Ausgabe aller zeichen in einer T1-Codierung (wichtig für die Ausgabe von Umlauten!)
\usepackage{graphicx}  			% Einbinden von Grafiken erlauben
%\usepackage{amsmath}
%\usepackage{amsfonts}
%\usepackage{amssymb}
\usepackage{mathpazo} 			% Einstellung der verwendeten Schriftarten
\usepackage{textcomp} 			% zum Einsatz von Eurozeichen u. a. Symbolen
\usepackage{listings}			% Datstellung von Quellcode mit den Umgebungen {lstlisting}, \lstinline und \lstinputlisting
\usepackage{xcolor} 			% einfache Verwendung von Farben in nahezu allen Farbmodellen
\usepackage[intoc]{nomencl} 	% zur Erstellung des Abkürzungsberzeichnisses
\usepackage{fancyhdr}			% Zusatzpaket zur Gestaltung von Fuß und Kopfzeilen

% -----------------------------------------------------------------------------------------------------------------
% Zum Aktualisieren des Abkürzungsverzeichnisses bitte auf der Kommandozeile folgenden Befehl aufrufen :
%  makeindex Studienarbeit.nlo -s nomencl.ist -o Studienarbeit.nls
% -----------------------------------------------------------------------------------------------------------------

% Hier die persönlichen Daten eingeben:

\newcommand{\titel}{BYOD}
\newcommand{\untertitel}{Konzeptionierung einer Entscheidungsempfehlung für ein mittelständiges Unternehmen}
\newcommand{\arbeit}{Studienarbeit}
\newcommand{\studiengang}{Informationstechnik}
\newcommand{\autor}{Nicolas Konle, Luka Kröger}
\newcommand{\matrikelnr}{MATRIKELNUMMERN}
\newcommand{\kurs}{TINF15B3}
\newcommand{\abgabe}{\today}
\newcommand{\betreuerdhbw}{Ralf Brune}

\newcommand{\jahr}{2018}			% für Angabe im Copyright-Vermerk der Titelseite

% Abkürzungen
\newcommand{\ua}{\mbox{u.\,a.\ }}
\newcommand{\zB}{\mbox{z.\,B.\ }}
\newcommand{\bs}{$\backslash$}

\renewcommand{\nomname}{Abkürzungsverzeichnis}

% -------------------------------------------------------------------------------------------
% Definition der Kopf- und Fußzeilen
\lhead{}								% Kopf links
\chead{}								% Kopf mitte
\rhead{\sffamily{\titel}}				% Kopf rechts
\lfoot{}								% Fuß links
\cfoot{\sffamily{\thepage}}				% Fuß mitte
\rfoot{\sffamily{\autor}}				% Fuß rechts
\renewcommand{\headrulewidth}{0.4pt}	% Liniendicke Kopf
\renewcommand{\footrulewidth}{0.4pt}	% Liniendicke Fuß


\makenomenclature							% Abkürzungsverzeichnis erstellen

% alle Abkürzungen, die in der Studienarbeit verwendet werden

\nomenclature{DHBW}{Duale Hochschule Baden-Württemberg}
\nomenclature{Sem}{Semester}
\nomenclature{OSS}{Open Source Software}					% Datei mit Abkürzungen laden

% -------------------------------------------------------------------------------------------
%                     Beginn des Dokumenteninhalts
% -------------------------------------------------------------------------------------------
\begin{document}
\setcounter{secnumdepth}{3}					% Nummerierungstiefe fürs Inhaltsverzeichnis
\setcounter{tocdepth}{3}
\sffamily									% für die Titelei serifenlose Schrift verwenden

% ------------------------------ Titelei -----------------------------------------------------

\thispagestyle{plain}
\begin{titlepage}
\enlargethispage{4.0cm}
\sffamily 								% Serifenlose Grundschrift für die Titelseite einstellen
				
\begin{flushright}
\includegraphics[scale=2.0]{Bilder/logo_dhbw.jpg}\\[5ex]
\end{flushright}

\begin{center}

\huge{\textsc{\textbf{\titel}}}\\[1.5ex]
\Large{\textbf{\untertitel}}\\[5ex]
\LARGE{\textbf{\arbeit}}\\[2ex]
\normalsize{für die Prüfung zum\\[1ex] Bachelor of Engineering}\\[3ex]
\Large{Studiengang \studiengang}\\[1ex]
\normalsize{Duale Hochschule Baden-Württemberg Karlsruhe}\\[5ex]
von\\[1ex] \autor \\[18ex]


\end{center}

\begin{flushleft}

\begin{tabular}{ll}
Abgabedatum:					& \quad \abgabe \\
Bearbeitungszeitraum:			& \quad 12 Wochen   \\ 
Matrikelnummer, Kurs: 			& \quad \matrikelnr , \kurs \\ 
Betreuer der Dualen Hochschule: & \quad \betreuerdhbw \\ [5ex]

\end{tabular} 



\small
Copyrightvermerk:\\

Dieses Werk einschließlich seiner Teile ist \textbf{urheberrechtlich geschützt}. Jede Verwertung außerhalb der engen Grenzen des Urheberrechtgesetzes ist ohne Zustimmung des Autors unzulässig und strafbar. Das gilt insbesondere für Vervielfältigungen, Übersetzungen, Mikroverfilmungen sowie die Einspeicherung und Verarbeitung in elektronischen Systemen.
\end{flushleft}
\begin{flushright}
\copyright{} \jahr
\end{flushright}
\end{titlepage} 				% erzeugt die Titelseite
\pagenumbering{Roman}						% große, römische Seitenzahlen für Titelei
\addchap{Eidesstattliche Erklärung}
Ich versichere hiermit, dass ich meine Studienarbeit mit dem Thema
\begin{quote}
\textit{\titel} -\textit{ \untertitel }
\end{quote}
selbständig verfasst und keine anderen als die angegebenen Quellen und Hilfsmittel benutzt habe. Die Arbeit wurde bisher keiner anderen Prüfungsbehörde vorgelegt und auch nicht veröffentlicht.


Mir ist bekannt, dass ich meine Diplomarbeit zusammen mit dieser Erklärung fristgemäß nach Vergabe des Themas in dreifacher Ausfertigung und gebunden im Sekretariat meines Studiengangs an der DHBW Karlsruhe abzugeben habe. Als Abgabetermin giltbei postalischer Übersendung der Eingangsstempel der DHBW, also nicht der Poststempel oder der Zeitpunkt eines
Einwurfs in einen Briefkasten der DHBW.\\[10ex]

Karlsruhe, den \today \\[4ex]


\rule[-0.2cm]{5cm}{0.5pt} \\

\textsc{\autor} \\[10ex]

% Sperrvermerk bei Bedarf dekommentieren
\hrule 
\vspace*{1.0cm}
\noindent \textbf{\Large{Sperrvermerk}}\\
\normalsize 				% Einbinden der eidestattlichen Erklärung
\chapter*{Abstract/Zusammenfassung} %*-Variante sorgt dafür, das Abstract nicht im Inhaltsverzeichnis auftaucht

Hier bitte den Abstract Ihrer Arbeit eintragen. Der Abstract sollte nicht länger als eine halbe Seite sein. Bitte klären Sie mit Ihrem Studiengangsleiter ab, ob der Abstract in englischer oder deutscher Sprache (oder möglicherweise sogar in beiden Sprachen) verfasst werden soll.
   				% Einbinden des Abstracts

\tableofcontents							% Erzeugen des Inhalsverzeichnisses
\printnomenclature[2.0cm]					% Erzeugen des Abkürzungsverzeichnisses
\listoffigures 								% Erzeugen des Abbildungsverzeichnisses 
\listoftables 								% Erzeugen des Tabellenverzeichnisses
\pagebreak

% --------------------------------------------------------------------------------------------
%                    Inhalt der Studienarbeit
%---------------------------------------------------------------------------------------------
\pagenumbering{arabic}						% arabische Seitenzahlen für den Hauptteil
\pagestyle{fancy}					
\rmfamily

\chapter{Einleitung}
\label{cha:Einleitung}

\section{Motivation}
\label{sec:Motivation}

\section{Ziel der Arbeit}
\label{sec:ZielDerArbeit}

\section{Aufbau der Arbeit}
\label{sec:AufbauDerArbeit}




\chapter{Vorlagen}
\label{cha:Vorlagen}

\section{Standards}
\subsection{Listenumgebungen und Fußnoten}
Jede wissenschaftliche Arbeit ist natürlich auf Fußnoten\footnote{das sind die kleinen zusätzlichen Hinweise am unteren Rand der Seite} angewiesen. Zudem kommt es immer wieder vor, dass man \marginpar{Bemerkung!}
\begin{itemize}
\item[-] Aufzählungen
\item[+] Nummerierungen oder
\item[*] Definitionen 
\end{itemize}
verwenden muss. In einer Aufzählung \footnote{also in einer \textit{enumerate}-Umgebung} würde das dann so aussehen.
\begin{enumerate}
\item Aufzählungen
\item Nummerierungen oder
\item Definitionen 
\end{enumerate}

In einer Definition \footnote{also in einer \textit{description}-Umgebung} sähe das dann wohl eher so aus:

\begin{description}
\item[Silvester]Jahresendfeier mit Feuerwerk und Alkoholgenuss
\item[Böller] Fuerwerkszubehör ohne visuellen Reiz, dafür aber recht laut
\end{description}
\subsection{Verweise und Zitate}
Natürlich muss man hin und wieder auch auf andere Kapitel verweisen so \zB in diesem Fall auf das Kapitel \ref{cha:wortberge} auf Seite \pageref{cha:wortberge}. Dazu muss das entsprechende Kapitel zuvor entsprechend mit dem Befehl \textit{\bs label\{Labelbezeichner\}} versehen worden sein. In \cite{foobar2003} wird dieser Fall bis ins kleinste Detail beschrieben.

% --------------------------------------------------------------------------------------------------------------------------

\section{Verschiedene Umgebungen}
\label{sec:Umgebungen}

\subsection{Einsatz von Programmlistings}
Für die Vorlage wird das paket \textit{listings} verwendet. \\

\begin{lstlisting} [language=PHP, numbers=left, numberstyle=\tiny, numbersep=10pt]
define('PATH_site', dirname(PATH_thisScript).'/');

if (@is_dir(PATH_site.'typo3/sysext/cms/tslib/')) {
        define('PATH_tslib', PATH_site.'typo3/sysext/cms/tslib/');
} elseif (@is_dir(PATH_site.'tslib/')) {
        define('PATH_tslib', PATH_site.'tslib/');
} else {
      
}
\end{lstlisting}

Das Paket \textit{listings} bietet zahlreiche Konfigurationsmöglichkeiten, um die Quellcodedarstellung an die eigenen Wünsche anzupassen. In einer fertig konfigurierten TexLive-Umgebung erfahren Sie mit dem Kommando

\begin{verbatim}
user@client:~> texdoc listings
\end{verbatim}

mehr über die Möglichkeiten des Pakets.

\subsection{Einsatz von Gleitumgebungen}
\subsubsection{Tabellen}

Tabellen selbst werden in der Umgebung \textit{tabular} oder \textit{tabularx}gesetzt. Um die Tabelle zu einem Gleitobjekt zu machen, muss diese dann in die Umgebung \textit{table} gesetzt werden. 

\begin{table}[hbt]
\centering
\begin{tabular}{c|c|c}
\hline Diese & Tabelle & ist \\ 
\hline zentriert & und  & verwendet \\ 
\hline vertikale & Trennzeichen &  .\\ 
\hline 
\end{tabular}
\caption{Beispiel für eine Tabelle} 

\end{table}


\subsubsection{Bilder}

Bilder werden mit dem Befehl \textit{\bs includecraphics} eingebunden. Um ein Bild zu einem Gleitobjekt zu machen, muss es in die Umgebung figure gesetzt werden.

\begin{figure}[hbt]
\centering
\includegraphics[scale=2.0]{Bilder/logo_dhbw}
\caption{Das Logo der DHBW Karlsruhe}
\end{figure}

Weit hinten, hinter den Wortbergen, fern der Länder Vokalien und Konsonantien leben die Blindtexte. Abgeschieden wohnen Sie in Buchstabhausen an der Küste des Semantik, eines grossen Sprachozeans. Ein kleines Bächlein namens Duden fliesst durch ihren Ort und versorgt sie mit den nötigen Regelialien.

Es ist ein paradiesmatisches Land, in dem einem gebratene Satzteile in den Mund fliegen. Nicht einmal von der allmächtigen Interpunktion werden die Blindtexte beherrscht - ein geradezu unorthographisches Leben.

Eines Tages aber beschloss eine kleine Zeile Blindtext, ihr Name war Lorem Ipsum, hinaus zu gehen in die weite Grammatik. Der grosse Oxmox riet ihr davon ab, da es dort wimmele von bösen Kommata, wilden Fragezeichen und hinterhältigen Semikoli, doch das Blindtextchen liess sich nicht beirren. Es packte seine sieben Versalien, schob sich sein Initial in den Gürtel und machte sich auf den Weg.

Als es die ersten Hügel des Kursivgebirges erklommen hatte, warf es einen letzten Blick zurück auf die Skyline seiner Heimatstadt Buchstabhausen, die Headline von Alphabetdorf und die Subline seiner eigenen Strasse, der Zeilengasse. Wehmütig lief ihm eine rethorische Frage über die Wange, dann setzte es seinen Weg fort.

Unterwegs traf es eine Copy. Die Copy warnte das Blindtextchen, da, wo sie herkäme wäre sie zigmal umgeschrieben worden und alles, was von ihrem Ursprung noch übrig wäre, sei das Wort "und" und das Blindtextchen solle umkehren und wieder in sein eigenes, sicheres Land zurückkehren.




\chapter{Weit hinter den Wortbergen}
\label{cha:wortberge}


Weit hinten, hinter den Wortbergen, fern der Länder Vokalien und Konsonantien leben die Blindtexte. Abgeschieden wohnen Sie in Buchstabhausen an der Küste des Semantik, eines grossen Sprachozeans. Ein kleines Bächlein namens Duden fliesst durch ihren Ort und versorgt sie mit den nötigen Regelialien.

Es ist ein paradiesmatisches Land, in dem einem gebratene Satzteile in den Mund fliegen. Nicht einmal von der allmächtigen Interpunktion werden die Blindtexte beherrscht - ein geradezu unorthographisches Leben.

Eines Tages aber beschloss eine kleine Zeile Blindtext, ihr Name war Lorem Ipsum, hinaus zu gehen in die weite Grammatik. Der grosse Oxmox riet ihr davon ab, da es dort wimmele von bösen Kommata, wilden Fragezeichen und hinterhältigen Semikoli, doch das Blindtextchen liess sich nicht beirren. Es packte seine sieben Versalien, schob sich sein Initial in den Gürtel und machte sich auf den Weg.

Als es die ersten Hügel des Kursivgebirges erklommen hatte, warf es einen letzten Blick zurück auf die Skyline seiner Heimatstadt Buchstabhausen, die Headline von Alphabetdorf und die Subline seiner eigenen Strasse, der Zeilengasse. Wehmütig lief ihm eine rethorische Frage über die Wange, dann setzte es seinen Weg fort.

Unterwegs traf es eine Copy. Die Copy warnte das Blindtextchen, da, wo sie herkäme wäre sie zigmal umgeschrieben worden und alles, was von ihrem Ursprung noch übrig wäre, sei das Wort "und" und das Blindtextchen solle umkehren und wieder in sein eigenes, sicheres Land zurückkehren.

Doch alles Gutzureden konnte es nicht überzeugen und so dauerte es nicht lange, bis ihm ein paar heimtückische Werbetexter auflauerten, es mit Longe und Parole betrunken machten und es dann in ihre Agentur schleppten, wo sie es für ihre Projekte wieder und wieder missbrauchten.

Und wenn es nicht umgeschrieben wurde, dann benutzen Sie es immer noch.


\textbf{}
%\begin{figure}[htb]
%\centering
%\includegraphics[width=0.8\textwidth]{FHWTLogo.jpg}
%\caption{Das Logo der FHWT}
%\label{fig:FHWTLogo}
%\end{figure}


\chapter{Zusammenfassung}

\section{}

\section{}



% ---------------------------- Literaturverzeichnis ----------------------------------------------

\begin{thebibliography}{999999}

\bibitem[FoBa03]{foobar2003} Foo, John; Bar, Belinda: \emph{Titel : Untertitel},\\ Verlagsort: Verlag, Jahr der Auflage. S. 10-20
\bibitem[Le01]{levy2001} Autor Name: \emph{Titel des Buches}, New York: Penguin Books, 2001

\end{thebibliography}

% ------------------------------- Anhang ---------------------------------------------------------

\begin{appendix}
\clearpage
\pagenumbering{Roman}						% römische Seitenzahlen für Anhang
\end{appendix}


\end{document}
